\documentclass{article}
\usepackage{amsmath, amssymb, amsthm}
\usepackage{hyperref}
\usepackage{geometry}
\geometry{margin=1in}

\title{Erd\H{o}s Problem \#313}
\date{Last updated: 2025-08-31}
\author{Source: erdosproblems.com}

\begin{document}
\maketitle

\noindent\textbf{Status:} OPEN \quad \textbf{No prize}

\noindent\textbf{Tags:} number theory, unit fractions

\medskip
\noindent\textbf{Problem Statement:}

Are there infinitely many solutions to\[\frac{1}{p_1}+\cdots+\frac{1}{p_k}=1-\frac{1}{m},\]where $m\geq 2$ is an integer and $p_1<\cdots<p_k$ are distinct primes?

\bigskip
\noindent\textbf{Remarks:}

For example,\[\frac{1}{2}+\frac{1}{3}=1-\frac{1}{6}\]and\[\frac{1}{2}+\frac{1}{3}+\frac{1}{7}=1-\frac{1}{42}.\]It is clear that we must have $m=p_1\cdots p_k$, and hence in particular there is at most one solution for each $m$. The integers $m$ for which there is such a solution are known as primary pseudoperfect numbers, and there are $8$ known, listed in A054377 at the OEIS.

\bigskip
\noindent\small{Source: \url{https://www.erdosproblems.com/313}}
\end{document}
